\scanpage{91}
\uline{NB} Si on a un hom
\[
u:B\longrightarrow\underline{\mathrm{Pic}}_{X/S}
\]
avec $B$ un schéma abélien, et si $\underline{\Omega}^{1}_{X/S}$ coh. plat $/S$ en dim $0$, alors $u$ se factorise en $B\longrightarrow\underline{\mathrm{Pic}}_{X/S}^{\omega}$. En particulier, si $\underline{\mathrm{Pic}}_{X/S}^{00}$ existe, alors on a
\[
\underline{\mathrm{Pic}}_{X/S}^{00}\subset\underline{\mathrm{Pic}}_{X/S}^{\omega}.
\]

Notons que la \uline{donnée} de $u$ équivaut : \margin{sous réserve d'existence de $B^{*}$ et de possibilité de descente (p.ex. si $B/S$ est projectif)} la donnée d'un torseur $A$ sous \deleted{$A$} le schéma dual $B^{*}$ de $B$, et d'un morphisme
\[
X\longrightarrow A
\]
\[
[\,\text{d'où}\qquad B\simeq\underline{\mathrm{Pic}}^{0}_{A/S}\subset\underline{\mathrm{Pic}}_{A/S}\longrightarrow\underline{\mathrm{Pic}}_{X/S}\quad\ldots\,]
\]
On trouve alors, par fonctorialité\margin{fonctorialité doit être explicitée avec un peu de soin !} un diagramme
\diagram{\begin{tikzcd}[ampersand replacement=\&,column sep=small,row sep=small]
0\arrow[r,dash] \& \underline{\omega}^{\unsure{*}}_{X/S}\arrow[r] \& \mathcal M_{X/S}\arrow[r] \& \underline{\mathrm{Pic}}^{\omega}_{X/S}\arrow[r] \& 0\\
0\arrow[r] \& \underline{\omega}_{A/S}\arrow[u]\arrow[r] \& \mathcal M_{A/S}\arrow[u]\arrow[r] \& \underline{\mathrm{Pic}}^{\omega}_{A/S}\arrow[u]\arrow[r,dash] \& 0\\
0\arrow[r,dash] \& \underline{\omega}_{A/S}\arrow[u,equal]\arrow[r] \& \deleted{\unsure{M_{B/S}|\underline E}}\mathcal M^{0}_{A/S}\arrow[u]\arrow[r] \& B\arrow[u,hook']\arrow[r] \& 0\\
\& V(R^{\unsure{1}}f_{B*}(\underline{\mathcal{O}}_{B}))\arrow[u,"\simeq" description,phantom] \& \& \underline{\mathrm{Pic}}^{0}_{A/S}\arrow[u,equal]
\end{tikzcd}}

Ici $\deleted{\unsure{M_{B/S}|\underline E_{B^*/S}}}\mathcal M^{0}_{A/S}$ est défini comme image inverse de $\mathcal M_{A/S}$ \inserted{par $B\hookrightarrow\underline{\mathrm{Pic}}^{\omega}_{A/S}$} [ou aussi de $\mathcal M_{X/S}$ par $B\longrightarrow\underline{\mathrm{Pic}}^{\omega}_{X/S}$ \inserted{si $\underline{\omega}_{A/S}\simeq\underline{\omega}_{X/S}$}]. Je vou-

% archive sheet number: 90

